\section{Conclusion}

We introduced FinIF, a benchmark suite and data construction framework for evaluating whether LLMs can produce usable financial work products under explicit delivery constraints.
FinIF-Test provides a 300-item hard benchmark with 3,134 instruction-following constraints, while FinIF-Train provides 1,064 workflow-grounded examples for training and expansion.
Experiments show that strong models satisfy most individual constraints but frequently fail strict item-level compliance.
The largest weaknesses involve evidence placement, quantitative verification, and decision-boundary reasoning.
These findings suggest that financial LLM evaluation should measure not only answer quality, but also whether the answer follows the operational instructions needed for auditability and downstream use.

\section*{Limitations}

FinIF has several limitations.
First, semantic evaluation depends on an LLM judge, so further human calibration and judge-sensitivity analysis are needed before treating scores as definitive.
Second, some private financial records are synthetic, which protects privacy and improves controllability but may not capture all institution-specific policies and data distributions.
Third, the current reported experiments cover five completed model runs; broader evaluation over open-weight, finance-tuned, multilingual, and newly released models should be added.
Fourth, FinIF-Test emphasizes hard instruction-following cases, so its strict pass rates should not be interpreted as general financial task success rates.
Finally, same-split supervised fine-tuning experiments on the new 300-item test split remain future work for the next version of this draft.

\section*{Ethical Considerations}

FinIF is designed for evaluating and improving instruction adherence in financial workflows, not for providing investment, legal, tax, or compliance advice.
The benchmark should therefore be used as an offline evaluation and research resource rather than as an automated decision system.
Synthetic private records are used to avoid exposing real customer or institutional data.
For public-source contexts, users of the dataset should respect the licenses and terms of the underlying materials.
Because higher instruction-following ability could make financial outputs appear more authoritative, deployed systems should still include human review, source verification, and institution-specific compliance controls.
